% =============================================
% Section 2: Related Work
% =============================================
\section{Related Work}
\label{sec:related}

\subsection{Small Language Models}

The question of how small a language model can be while retaining meaningful
generation capability has received growing attention. Eldan and
Li~\cite{eldan2023tinystories} introduced the TinyStories dataset---a
corpus of short stories using vocabulary accessible to 3--4 year olds,
generated by GPT-3.5 and GPT-4. They demonstrated that models with as few
as 2.5M parameters can produce fluent, grammatically correct stories when
trained on this focused dataset. Their evaluation framework, which uses
GPT-4 as a multi-dimensional judge (assessing grammar, creativity, and
consistency), established new benchmarks for evaluating sub-10M parameter models.

The Phi series of models~\cite{gunasekar2023textbooks} extended this line of
work, showing that carefully curated ``textbook-quality'' training data enables
small models to punch above their weight class. These results collectively
suggest that the coherence floor for language models may be substantially lower
than previously assumed, provided the training data and evaluation methodology
are appropriately designed.

\subsection{Low-Bit and Ternary LLM Architectures}

The BitNet architecture~\cite{wang2023bitnet} introduced 1-bit weight
quantization for Transformers, demonstrating that binary weights could scale
competitively with full-precision counterparts. BitNet
b1.58~\cite{ma2024bitnet} refined this approach by adopting ternary weights
$\{-1, 0, +1\}$, achieving 1.58 bits per parameter. This ternary scheme
preserves the ability to represent zero (enabling sparse computation) while
maintaining model quality competitive with FP16 LLMs of equivalent size.
The recently released BitNet b1.58 2B4T~\cite{bitnet2b4t2025} provides the
first open-source native 1-bit LLM at 2B parameters, trained on 4 trillion
tokens.

Zhu et al.~\cite{zhu2024matmulfree} took a more radical approach with their
Matmul-free LLM, entirely eliminating matrix multiplication operations from
the Transformer architecture. They replaced dense layers with ternary-weight
additive operations and self-attention with Hadamard (element-wise) products,
achieving up to 61\% memory reduction during training and over 10$\times$
reduction during inference, while maintaining competitive perplexity.

These architectures are particularly relevant for microcontroller deployment,
as they replace the most expensive operations (floating-point matrix
multiplication) with operations that map naturally to the integer and bitwise
units available on embedded processors.

\subsection{Language Model Inference on Microcontrollers}

Karpathy's \texttt{llama2.c}~\cite{karpathy2023llama2c} provided a seminal
reference implementation: a pure-C inference engine for LLaMA-architecture
models in under 500 lines. While designed primarily for educational purposes,
it demonstrated that Transformer inference could be stripped to its essential
operations. Community efforts have since ported \texttt{llama2.c} to various
embedded platforms, including the \espss{}, though typically limited to the
smallest (260K parameter) model.

MambaLite-Micro~\cite{mambalite2024} represents the state of the art in MCU
language model deployment. By implementing a runtime-free Mamba (state space
model) inference engine in pure C, they achieved 83\% peak memory reduction
compared to Transformer baselines on both \espss{} and STM32H7 platforms.
Their key insight was eliminating large intermediate tensors through fused operators and memory layout optimization. While highly efficient, these models typically target classification tasks such as Keyword Spotting (KWS) or activity recognition, which are fundamentally different from the open-ended generative reasoning targeted by \tinybit{}.

\subsection{Edge LLMs vs. Legacy Keyword Spotting}

Modern voice-controlled appliances (e.g., smart fans, air conditioners) predominantly rely on Keyword Spotting (KWS), which uses fixed audio template matching for a predefined set of commands (e.g., ``Power On''). While power-efficient, KWS lacks semantic flexibility; it cannot process natural language variations (e.g., ``I'm feeling a bit hot'') or complex multi-step instructions. In contrast, \tinybit{} deploys a general-purpose generative brain on the \espss{}, enabling not only zero-shot intent understanding but also cross-domain reasoning and even instruction-following for device maintenance, which provides a significant scalability advantage over rigid KWS implementations.

The \texttt{llama4micro} project~\cite{llama4micro} and related community
efforts have explored running quantized LLMs on MCUs, but no prior work
has systematically studied the relationship between model capacity,
quantization level, and language coherence quality on microcontroller
hardware, nor has any work developed ISA-specific optimized kernels for
ternary inference on Xtensa processors.

\subsection{Positioning of This Work}

Table~\ref{tab:comparison} summarizes the positioning of \tinybit{} relative
to prior work. Our contribution uniquely combines (1)~a systematic coherence
analysis for sub-10M models under quantization, with (2)~ISA-native kernel
optimization for a specific MCU target, providing both the theoretical
understanding and practical toolchain needed for coherent language generation
on microcontrollers.

\begin{table*}[t]
\centering
\caption{Comparison of \tinybit{} with related work in efficient LLM inference.}
\label{tab:comparison}
\small
\begin{tabular}{lccccc}
\toprule
\textbf{Work} & \textbf{Target} & \textbf{Ternary} & \textbf{MCU Kernel} & \textbf{Coherence} & \textbf{Text Gen.} \\
& & \textbf{Weights} & \textbf{Optimization} & \textbf{Analysis} & \textbf{on MCU} \\
\midrule
TinyStories~\cite{eldan2023tinystories}    & GPU/CPU & \texttimes & \texttimes & \checkmark & \texttimes \\
BitNet b1.58~\cite{ma2024bitnet}           & GPU     & \checkmark & \texttimes & \texttimes & \texttimes \\
Matmul-free LLM~\cite{zhu2024matmulfree}   & GPU     & \checkmark & \texttimes & \texttimes & \texttimes \\
\texttt{llama2.c}~\cite{karpathy2023llama2c} & CPU/MCU & \texttimes & \texttimes & \texttimes & \checkmark$^*$ \\
MambaLite-Micro~\cite{mambalite2024}       & MCU     & \texttimes & \checkmark$^\dagger$ & \texttimes & \texttimes \\
\midrule
\textbf{\tinybit{} (ours)}                 & \textbf{MCU}  & \checkmark & \checkmark & \checkmark & \checkmark \\
\bottomrule
\multicolumn{6}{l}{\footnotesize $^*$Limited to 260K-parameter model. \quad $^\dagger$SSM-specific, not ternary kernel.}
\end{tabular}
\end{table*}
